\subsection{Economic Potential}

The economic potential assessment determines the financial viability of lake-based district heating systems compared to individual small-scale air-source heat pump installations at each building. This analysis builds upon the cost-benefit comparison methodology from \cite{Eggimann.2023}, extending it to incorporate annualized cost comparisons with differentiated financing conditions for decentralized and centralized systems. The assessment evaluates economic feasibility across 1-hectare raster cells within a 10-kilometer radius of each lake.

The economic viability is assessed by comparing the total annualized costs of decentralized air-source heat pump supply against centralized lake-source district heating supply. The annuity factor converts investment costs into equivalent annual payments and is defined as:

\begin{equation}
a = \frac{r \cdot (1 + r)^T}{(1 + r)^T - 1}
\label{eq:annuity}
\end{equation}

\noindent where $r$ is the weighted average cost of capital (WACC) and $T$ is the system lifetime. Separate annuity factors are applied for decentralized systems ($a_d$, using $r_d$) and centralized systems ($a_c$, using $r_c$), reflecting the different financing conditions of private homeowners and large-scale utility operators.

The thermal capacity required to satisfy the heat demand of a raster cell is derived from the annual heat demand and assumed full-load hours:

\begin{equation}
P = \frac{E_\mathrm{h}}{h}
\end{equation}

\noindent where $E_\mathrm{h}$ is the annual heat demand [kWh/yr] and $h$ the technology-specific full-load hours [h/yr]. The annualized capital expenditure for each technology is then:

\begin{equation}
C_\mathrm{CAPEX,ann} = a \cdot C_\mathrm{inv} \cdot \frac{E_\mathrm{h}}{h}
\end{equation}

The annual operational expenditure (electricity cost) for each system is:

\begin{equation}
C_\mathrm{OPEX} = C_\mathrm{el} \cdot \frac{E_\mathrm{h}}{COP}
\end{equation}

For the centralized water-source heat pump system, the heat demand includes distribution losses:

\begin{equation}
E_\mathrm{h,Water} = E_\mathrm{h} \cdot (1 + q_\mathrm{loss})
\end{equation}

The economic viability index (EI) compares the annualized total cost of decentralized supply to centralized supply, normalized by the annualized pipeline investment required to connect the demand location to the lake:

\begin{equation}
\mathrm{EI} = \frac{
\left( a_d \cdot C_\mathrm{inv,Air} \cdot \dfrac{E_\mathrm{h}}{h_\mathrm{Air}} + C_\mathrm{el,Air} \cdot \dfrac{E_\mathrm{h}}{COP_\mathrm{Air}} \right)
- \left( a_c \cdot C_\mathrm{inv,Water} \cdot \dfrac{E_\mathrm{h,Water}}{h_\mathrm{Water}} + C_\mathrm{el,Water} \cdot \dfrac{E_\mathrm{h,Water}}{COP_\mathrm{Water}} \right)

}{a_c \cdot C_\mathrm{pipe} \cdot d}
\label{eq:EI}
\end{equation}

\noindent where $C_\mathrm{inv,Air}$ and $C_\mathrm{inv,Water}$ are specific investment costs [€/kW\textsubscript{th}], $h_\mathrm{Air}$ and $h_\mathrm{Water}$ are full-load hours [h/yr], $C_\mathrm{el,Air}$ and $C_\mathrm{el,Water}$ are electricity prices [€/kWh], $COP_\mathrm{Air}$ and $COP_\mathrm{Water}$ are coefficients of performance, $C_\mathrm{pipe}$ is the pipeline construction cost per unit length [€/m], and $d$ is the effective distance from the raster cell to the nearest lake [m]. The numerator represents the annual cost saving of centralized over decentralized heat supply [€/yr], while the denominator represents the annualized infrastructure investment required for the pipeline connection [€/yr]. Areas where the EI exceeds unity indicate that the annual cost savings from the centralized water-source system exceed the annualized pipeline costs, rendering the lake-based solution economically superior.

The use of separate financing rates reflects the structural difference between private homeowner investments (lower WACC due to subsidized lending and lower risk perception) and large-scale utility infrastructure projects (higher WACC reflecting capital market requirements and project risk). This differentiation provides a more realistic comparison than a uniform discount rate assumption.

The specific parameter values used in this German assessment are specified for German market conditions and assume a district heating network operating at approximately 75--80\,°C, which can supply both space heating and domestic hot water as is typical for most modern district heating networks in Germany utilizing large-scale water-source heat pumps \citep{Ahrendts.2022}. The methodology could also be applied to other technical configurations by adapting the COP values, costs for large-scale heat pumps, pipeline specifications, and distribution losses. The parameter values used in this assessment are detailed in Table~\ref{tab:economic_parameters}.

\begin{table}[htb]
\centering
\caption{Economic assessment parameter values for Germany-specific analysis}
\label{tab:economic_parameters}
\begin{tabular}{@{}lll@{}}
\toprule
Parameter & Value & Assumptions based on \\
\midrule
\multicolumn{3}{l}{\textit{Air-source heat pump (decentralized)}} \\
$C_\mathrm{el,Air}$ & 0.37\,€/kWh & \cite{BDEW.2026} \\
$C_\mathrm{inv,Air}$ & 1{,}872\,€/kW\textsubscript{th} & \cite{Thomsen.2022, dena.2025} \\
$h_\mathrm{Air}$ & 2{,}200\,h/a & \cite{Thomsen.2022} \\
$COP_\mathrm{Air}$ & 2.9 & \cite{Thomsen.2022, Meyer.2024, Olympios.2022} \\
$r_d$ (WACC decentral) & 2.2\,\% & \cite{dena.2025} \\
\addlinespace
\multicolumn{3}{l}{\textit{Water-source heat pump (centralized)}} \\
$C_\mathrm{el,Water}$ & 0.16\,€/kWh & \cite{BDEW_2.2026} \\
$C_\mathrm{inv,Water}$ & 1{,}064\,€/kW\textsubscript{th} & \cite{Thomsen.2022, dena.2025} \\
$h_\mathrm{Water}$ & 3{,}500\,h/a & \cite{Ahrendts.2022} \\
$COP_\mathrm{Water}$ & 3.3 & Own calculations based on Equation~\ref{eq:COP} \\
$r_c$ (WACC central) & 6.0\,\% & \cite{dena.2025} \\
\addlinespace
\multicolumn{3}{l}{\textit{Infrastructure and system}} \\
$C_\mathrm{pipe}$ & 2{,}500\,€/m & \cite{dena.2025} \\
$q_\mathrm{loss}$ & 12\,\% & \cite{dena.2025} \\
$T$ & 20\,years & \cite{dena.2025} \\
\bottomrule
\end{tabular}
\end{table}

The final economic potential for each lake represents the spatial aggregation of demand areas where lake-based heating demonstrates economic superiority, constrained by the sustainable thermal capacity determined in the technological assessment:

\begin{equation}
Q_{\text{economic}} = \sum_{i=1}^{N} \min\left(Q_{\text{viable}, i},\, Q_{\text{sustainable}, i}\right)
\end{equation}

\noindent where $Q_{\text{viable}, i}$ represents the heat demand in raster cell $i$ where $\mathrm{EI} > 1$, and $Q_{\text{sustainable}, i}$ represents the sustainable extraction capacity allocated to that cell. This approach ensures that economic potential remains within environmental sustainability boundaries while identifying realistic deployment opportunities based on spatially explicit cost-benefit analysis.

The economic viability assessment employs two distinct methodologies for calculating the distance $d$ between heat demand locations and lake thermal resources. The lake-distance approach calculates direct Euclidean distances from each raster cell to the nearest lake boundary. This method assumes individual pipeline connections from each demand location to the closest lake shore, providing a conservative estimate of distribution costs. While computationally efficient and spatially explicit, this approach may overestimate infrastructure requirements by not accounting for potential economies of scale through shared distribution networks.

The hub-based methodology introduces intermediate collection points that aggregate heat demand before connecting to lake thermal resources, more accurately reflecting realistic district heating network topologies. Heat demand pixels are grouped into spatial clusters (hubs) using a regular grid approach with configurable resolution (default: 1\,km\,$\times$\,1\,km grid cells). For each grid cell containing positive heat demand, a hub location is determined as the heat-demand-weighted centroid of all demand pixels within that cell.

The total effective pipeline distance for each pixel $i$ is calculated as:
\begin{equation}
d_{i} = d_{i,\mathrm{hub}} + s_i \cdot d_{\mathrm{hub,lake}}
\end{equation}

\noindent where $d_{i,\mathrm{hub}}$ represents the Euclidean connection distance from the demand pixel to its assigned hub, $d_{\mathrm{hub,lake}}$ is the distance from the hub to the nearest lake shore, and $s_i$ is the cost-sharing factor for pixel $i$. The sharing factor is calculated as $s_i = E_{h,i} / E_{h,\mathrm{hub}}$, where $E_{h,\mathrm{hub}}$ is the total heat demand of all pixels assigned to the same hub. This weighted sharing reflects that pixels with higher demand bear a proportionally larger share of the hub-to-lake pipeline cost, analogous to demand-proportional infrastructure cost allocation in district heating networks.